% Apéndice A — cómo se produjo este informe (se incluye desde informe.tex)
Este informe se produjo con el mismo tipo de herramientas que analiza, y el proceso tuvo sus propias
fallas; las registramos por coherencia.

\paragraph{Organización.} Una sesión de Claude Code (Opus~5.5), aislada en una rama propia y sin permiso
para hacer \emph{commits} sin autorización explícita, actuó como coordinadora. Primero convirtió las
transcripciones locales (unos 60\,MB de JSONL) en texto legible (3\,MB) con un script propio, excluido
del control de versiones porque las transcripciones pueden contener rutas internas del servidor de la
colaboración. Luego escribió una consigna común (\code{notes/AGENT\_BRIEF.md}) con el estándar de
evidencia y una plantilla de notas, y lanzó ocho agentes de lectura en paralelo:
\begin{center}\small
\begin{tabular}{ll}
\toprule
Agente & Fuentes\\
\midrule
T1, T2, T3 & transcripciones de Claude (29/08--02/09, 03/09--10/09, 10/09--23/09)\\
C1, C2 & transcripciones de Codex (hilo principal de Astra; sesiones posteriores)\\
R & artefactos en \code{claude\_work/} (linaje de cada afirmación física)\\
G & \code{git}, \emph{pull requests}, \code{CLAUDE.md}, \code{AGENTS.md}, memoria, prompts\\
P & cuaderno de laboratorio, evolución de scripts, notas GAP y capítulos de la tesis\\
\bottomrule
\end{tabular}
\end{center}
Un noveno agente verificó contra las fuentes primarias 21 afirmaciones clave (20 confirmadas, 1 con
matices) y un décimo revisó el primer borrador en modo adversarial.

\paragraph{Fallas del propio proceso.}
\begin{itemize}
  \item \textbf{Nuestra consigna tenía el signo al revés.} El resumen inicial decía que \emph{quitar} el
  requisito producía la inversión; es al revés. Lo detectó el primer agente en terminar, contrastando con
  los números de las sesiones; corregimos la consigna y avisamos a los demás agentes en curso. Es el mismo
  tipo de error que estudia el informe: una afirmación del coordinador que, sin verificación, se habría
  propagado a todos los agentes.
  \item \textbf{Límite de uso.} Siete de los ocho agentes fueron interrumpidos por el límite de sesión de la
  cuenta; seis ya habían escrito sus notas completas y dos se relanzaron horas después con la instrucción
  de escribir la nota antes del resumen final.
  \item \textbf{Fricción de herramientas.} El aislamiento de la sesión rechazaba cualquier comando cuyo
  texto contuviera la cadena «git» (incluida la ruta \code{/Github/}); se resolvió con scripts auxiliares.
  \item \textbf{Una cita inventada por una herramienta.} El resumen automático de una página web nos atribuyó
  a Schwartz la frase «se atraparon errores mutuamente», que no figura en el artículo. La detectó el revisor
  adversarial y la eliminamos tras reverificar textualmente todas las citas del curso.
  \item \textbf{El primer borrador tenía errores de fondo.} La revisión adversarial (61 observaciones) encontró
  que el borrador atribuía el sesgo sólo a la línea \code{HasStation} (la selección estaba también en la
  estructura de los datos), omitía que todo es Monte Carlo y que el UMD también está seleccionado, recortaba
  una cita condicional, afirmaba una relación causal que la transcripción contradice y tenía tres conteos mal.
  Esta versión incorpora esas correcciones.
  \item \textbf{Fuente inaccesible.} El enlace compartido de ChatGPT devolvió un error de acceso; usamos el
  registro local de Codex; según el tesista todo el trabajo se hizo por terminal, así que es la misma
  conversación.
\end{itemize}
